% archon:future-covers MR4689373HigherSiegelWeilFormulaForUnitaryGroupsTheNonSingularTerms/PerverseSheaves.lean
\chapter{Perverse Sheaves \texorpdfstring{$\cK_d^{\mathrm{Eis}}$}{K-Eis} and \texorpdfstring{$\cK_d^{\mathrm{Int}}$}{K-Int}}
\label{chap:perverse-sheaves-K}

This chapter formalizes Phase~P3b of the project: the construction of two graded perverse
sheaves on \(\Herm_{2d}(X'/X)\), one encoding the analytic (Eisenstein) side of the formula
and one encoding the geometric (intersection) side. The key results are the Frobenius trace
formulas (1.2) and (1.3) from the paper, which translate the main theorem into an isomorphism
of perverse sheaves. We follow \S\S11--14 of Feng--Yun--Zhang [MR4689373].

\section{Hitchin-type moduli spaces}

\begin{definition}[Hitchin moduli stack]
  \label{def:hitchin-stack}
  \lean{HSW.HitchinStack}
  \uses{def:herm-torsion-stack, def:etale-cover, def:hermitian-shtukas-stack}
  % SOURCE: FYZ-HigherSiegelWeil.tex, §11.1–11.2
  % SOURCE QUOTE: "the Hitchin moduli stack \(\cM_d\) parameterizes pairs \((\cL,\phi)\)
  %   with \(\cL\) a line bundle on \(X'\) and
  %   \(\phi:\cL\to\sigma^*\cL^\vee\) an injective map with
  %   \(\coker(\phi)\in\Herm_{2d}(X'/X)\)"
  \textit{Source: Feng--Yun--Zhang [MR4689373], \S11.1--11.2.}
  For \(d \geq 0\), the \emph{Hitchin moduli stack} \(\cM_d\) parameterizes pairs
  \((\cL, \phi)\) where:
  \begin{enumerate}
    \item \(\cL\) is a line bundle on \(X'\);
    \item \(\phi : \cL \to \sigma^*\cL^\vee\) is an injective sheaf map (a rank-1
      Hermitian Higgs field) whose cokernel
      \(\cQ := \coker(\phi) \in \Herm_{2d}(X'/X)\)
      is a Hermitian torsion sheaf of length \(2d\).
  \end{enumerate}
  Here \(\sigma^*\cL^\vee = \sigma^*\underline{\Hom}(\cL, \om_{X'})\) uses the Serre
  dual on \(X'\). The stack \(\cM_d\) is an algebraic stack of finite type over
  \(\mathbf{F}_q\), fibered over \(\Herm_{2d}(X'/X)\) via the Hitchin map
  \(f_d : \cM_d \to \Herm_{2d}\) of Definition~\ref{def:hitchin-map}.
\end{definition}

\begin{definition}[Hitchin map]
  \label{def:hitchin-map}
  \lean{HSW.HitchinMap}
  \uses{def:hitchin-stack, def:herm-torsion-stack}
  % SOURCE: FYZ-HigherSiegelWeil.tex, §11.3
  % SOURCE QUOTE: "the Hitchin map \(f_d:\cM_d\to\Herm_{2d}(X'/X)\) sends
  %   \((\cL,\phi)\mapsto\coker(\phi)\)"
  \textit{Source: Feng--Yun--Zhang [MR4689373], \S11.3.}
  The \emph{Hitchin map} is the proper morphism of algebraic stacks
  \[
    f_d : \cM_d \to \Herm_{2d}(X'/X),
    \qquad (\cL, \phi) \mapsto \coker(\phi).
  \]
  Properness holds because the cokernel sheaf lies in a bounded family: fixing the
  degree of \(\cL\) and the length \(2d\) of \(\coker(\phi)\) gives a proper moduli
  problem.
\end{definition}

\section{Local intersection number and trace formula}

\begin{definition}[Intersection perverse sheaf \texorpdfstring{$\cK_d^{\mathrm{Int}}$}{K-Int}]
  \label{def:K-Int}
  \lean{HSW.KInt}
  \uses{def:hitchin-map, def:herm-torsion-stack, def:springer-sheaf-herm}
  % SOURCE: FYZ-HigherSiegelWeil.tex, §11.4, equation (1.3)
  % SOURCE QUOTE: "\(\cK_d^\Int=(f_d)_*\Ql\in\Perv(\Herm_{2d},\Ql)\)
  %   ... is a direct summand of \(\Spr^\Herm_{2d}\)"
  \textit{Source: Feng--Yun--Zhang [MR4689373], \S11.4, eq.~(1.3).}
  The \emph{intersection perverse sheaf} is
  \[
    \cK_d^{\mathrm{Int}} := (f_d)_*\, \Ql
    \;\in\; \mathrm{Perv}(\Herm_{2d}(X'/X),\, \Ql).
  \]
  By the Decomposition Theorem and the Hermitian analogue of the geometric Satake
  isomorphism, \(\cK_d^{\mathrm{Int}}\) is a direct summand of the Hermitian
  Springer sheaf \(\mathrm{Spr}^{\mathrm{Herm}}_{2d}\)
  (Definition~\ref{def:springer-sheaf-herm}) and carries the induced \(W_d\)-action.
\end{definition}

\begin{definition}[Eisenstein perverse sheaf \texorpdfstring{$\cK_d^{\mathrm{Eis}}$}{K-Eis}]
  \label{def:K-Eis}
  \lean{HSW.KEis}
  \uses{def:herm-torsion-stack, def:siegel-eisenstein-series, def:degenerate-principal-series,
        def:springer-sheaf-herm}
  % SOURCE: FYZ-HigherSiegelWeil.tex, Def. (def: K_Eis); Thm. (th:Den Wd); eq. (eq: comparison 1)
  % SOURCE QUOTE: "\(\wt{E}_{a}(m(\cE),s,\Phi) =  \Tr(\Frob,\cK_{d}^{\Eis}(q^{-2s})_{\cQ})\)"
  \textit{Source: Feng--Yun--Zhang [MR4689373], Def.~``K\_Eis''; Thm.~``Den Wd''; eq.~(comparison~1).}
  The \emph{Eisenstein perverse sheaf} is the unique perverse sheaf
  \[
    \cK_d^{\mathrm{Eis}} \in \mathrm{Perv}(\Herm_{2d}(X'/X),\, \Ql)
  \]
  that is a direct summand of \(\mathrm{Spr}^{\mathrm{Herm}}_{2d}\) and whose
  Frobenius trace at a \(k\)-point \(\cQ \in \Herm_{2d}(X'/X)(\mathbf{F}_q)\)
  encodes the \(d\)-th Fourier coefficient of the Siegel--Eisenstein series:
  \[
    \Tr(\Frob,\; (\cK_d^{\mathrm{Eis}}(q^{-2s}))_{\overline{\cQ}})
    \;=\;
    \widetilde{E}_a(m(\cE), s, \Phi)
  \]
  (a plain stalk trace, with no automorphism weight; here \(\cE = \cE(\cQ)\) is the
  rank-1 bundle whose cokernel of the Hermitian map is \(\cQ\), and
  \(\widetilde{E}_a\) is the normalized Fourier coefficient of
  Definition~\ref{def:siegel-eisenstein-series}).
  The Eisenstein sheaf is constructed via the geometric Langlands / Eisenstein functor
  applied to the degenerate principal series \(I_1(s, \chi)\) of
  Definition~\ref{def:degenerate-principal-series}.
\end{definition}

\begin{theorem}[Frobenius trace formula: Eisenstein side]
  \label{thm:trace-Eis}
  \lean{HSW.traceFrob_Eis}
  \uses{def:K-Eis, def:normalized-fourier-coeff, def:central-value-d,
        def:herm-torsion-stack, def:fourier-coefficient}
  % SOURCE: FYZ-HigherSiegelWeil.tex, §12.3, equation (1.2)
  % SOURCE QUOTE: "\(\wt E_a(m(\cE),s,\Phi)=\sum_{i=0}^d
  %   \Tr(\Frob_\cQ,(\cK^\Eis_{d,i})_\cQ)q^{-2is}\)"
  \textit{Source: Feng--Yun--Zhang [MR4689373], \S12.3, eq.~(1.2).}
  For \(\cQ \in \Herm_{2d}(\mathbf{F}_q)\) and \(d \geq 0\), the graded Frobenius
  trace formula holds:
  \[
    \widetilde{E}_a(m(\cE), s, \Phi)
    \;=\;
    \sum_{i=0}^{d}
    \Tr\!\bigl(\Frob_q,\; (\cK^{\mathrm{Eis}}_{d,i})_{\overline{\cQ}}\bigr)\, q^{-2is},
  \]
  where \(\cK^{\mathrm{Eis}}_d = \bigoplus_{i=0}^d \cK^{\mathrm{Eis}}_{d,i}\) is the
  graded decomposition of \(\cK_d^{\mathrm{Eis}}\) by the eigenvalue \(q^{-2i}\) of the
  \(q^{ds}\)-factor and \(d = d(\cE) = -\deg(\cE) + n\deg\om_X\)
  (Definition~\ref{def:central-value-d}).
\end{theorem}

\begin{proof}
  \uses{def:K-Eis, def:normalized-fourier-coeff}
  By definition of \(\cK_d^{\mathrm{Eis}}\): the Frobenius trace at \(\cQ\) computes
  the normalized Fourier coefficient \(\widetilde{E}_a\) by the function-sheaf
  dictionary (Grothendieck--Lefschetz trace formula for \(\ell\)-adic sheaves) as a
  plain stalk trace at \(\cQ\). The grading by \(q^{-2is}\) records
  the \(\ell\)-adic weight filtration, and the sum over graded pieces gives the
  stated formula.
\end{proof}

\begin{theorem}[Frobenius trace formula: intersection side]
  \label{thm:trace-Int}
  \lean{HSW.traceFrob_Int}
  \uses{def:K-Int, def:degree-special-cycle, def:virtual-class-full, def:herm-torsion-stack}
  % SOURCE: FYZ-HigherSiegelWeil.tex, §1.3, eq. (intro KInt); Cor. (cor: KR degree as trace of Frob)
  % SOURCE QUOTE: "\(\deg [\cZ_{\cE}^r(a)]=\sum_{i=0}^{d}\Tr(\Frob_{\cQ}, (\cK^{\Int}_{d,i})_{\cQ})\cdot (d-2i)^{r}\)"
  \textit{Source: Feng--Yun--Zhang [MR4689373], \S1.3, eq.~(intro KInt); Cor.~``KR degree as trace of Frob''.}
  For \(\cQ \in \Herm_{2d}(\mathbf{F}_q)\) and any \(r \geq 0\)
  (the formula is uniform in \(r\); the grading sum runs over \(i = 0, \ldots, d\)
  regardless of \(r\)), the graded Frobenius trace formula holds:
  \[
    \deg[\cZ_{\cE}^r(a)]
    \;=\;
    \sum_{i=0}^{d}
    \Tr\!\bigl(\Frob_q,\; (\cK^{\mathrm{Int}}_{d,i})_{\overline{\cQ}}\bigr)
    \cdot (d - 2i)^r,
  \]
  where \(\cK^{\mathrm{Int}}_d = \bigoplus_{i=0}^d \cK^{\mathrm{Int}}_{d,i}\) is graded
  by eigenvalues \((d-2i)\) of the canonical endomorphism \(\bar{C}\) on
  \(Rf_{d*}\Ql\) (the Hitchin endomorphism), and
  \(\deg[\cZ_{\cE}^r(a)]\) is the degree of the virtual fundamental class from
  Definition~\ref{def:virtual-class-full}.
\end{theorem}

\begin{proof}
  \uses{def:K-Int, def:virtual-class-full}
  By the Grothendieck--Lefschetz trace formula, the Frobenius trace of
  \((f_{d*}\Ql)_{\overline{\cQ}}\) counts \(\mathbf{F}_q\)-points of the fiber
  \(f_d^{-1}(\cQ)\) weighted by \(|\mathrm{Aut}|^{-1}\). By the virtual class
  construction (Definition~\ref{def:virtual-class-full}), this weighted count equals
  \(\deg[\cZ_{\cE}^r(a)]\). The grading by \((d-2i)^r\) arises from the Hitchin
  endomorphism \(\bar{C}\) on the Hitchin fibration: its eigenvalues on the fiber
  stalk are \((d-2i)\), and raising to the \(r\)-th power recovers the factor
  needed to match the \(r\)-th derivative at \(s = 0\).
\end{proof}
